
\documentclass[11pt,a4paper]{article}
\usepackage[utf8]{inputenc}
\usepackage[T1]{fontenc}
\usepackage[french]{babel}
\usepackage{lmodern}
\usepackage{microtype}
\usepackage[a4paper,margin=1.65cm,headheight=15pt]{geometry}
\usepackage{amsmath,amssymb,mathtools}
\usepackage{array,tabularx,multicol,booktabs}
\usepackage{enumitem}
\usepackage[most]{tcolorbox}
\usepackage{xcolor}
\usepackage{tikz}
\usetikzlibrary{arrows.meta,calc,positioning}
\usepackage{pdflscape}
\usepackage{fancyhdr}
\usepackage{hyperref}
\usepackage{pagecolor}
\usepackage{needspace}

\definecolor{fondpage}{RGB}{247,248,250}
\definecolor{encre}{RGB}{31,35,40}
\definecolor{bleu}{RGB}{40,91,135}
\definecolor{vert}{RGB}{55,125,92}
\definecolor{orange}{RGB}{178,105,42}
\definecolor{violet}{RGB}{101,77,145}
\definecolor{rouge}{RGB}{170,72,72}
\definecolor{gris}{RGB}{86,91,99}
\definecolor{bleufonce}{RGB}{28,55,120}
\definecolor{bleuclair}{RGB}{238,243,255}
\definecolor{grisclair}{RGB}{245,245,245}

\hypersetup{colorlinks=true,linkcolor=bleufonce,urlcolor=bleufonce,
pdftitle={Parcours de revision -- Variables aléatoires -- Premiere specialite},pdfauthor={}}

\setlength{\parindent}{0pt}
\setlength{\parskip}{0.25em}
\renewcommand{\arraystretch}{1.13}
\setlength{\tabcolsep}{4pt}
\setlist[itemize]{leftmargin=1.25em,itemsep=0.15em,topsep=0.2em}
\setlist[enumerate,1]{label=\textbf{\arabic*.}, leftmargin=1.05cm, itemsep=0.35em, topsep=0.2em}
\setlist[enumerate,2]{label=\textbf{\alph*.}, leftmargin=0.85cm, itemsep=0.25em, topsep=0.15em}

\pagestyle{fancy}
\fancyhf{}
\cfoot{\thepage}

\newcommand{\R}{\mathbb{R}}
\newcommand{\N}{\mathbb{N}}
\newcommand{\sourceq}[1]{ {\scriptsize\textcolor{bleufonce}{(site Premiere q#1)}}}
\newcounter{question}
\newcounter{reponse}

\newenvironment{blocquestions}[1]{%
  \begin{tcolorbox}[enhanced,breakable,colback=bleuclair,colframe=bleufonce,boxrule=0.6pt,arc=2mm,left=2mm,right=2mm,top=2mm,bottom=1.5mm,title={\bfseries #1},coltitle=white,colbacktitle=bleufonce,before skip=3mm,after skip=3mm, fontupper=\large]
  \large
  \begin{enumerate}[leftmargin=*,label=\textbf{\arabic*.},itemsep=0.82em,topsep=0.5em]
  \setcounter{enumi}{\value{question}}
}{%
  \setcounter{question}{\value{enumi}}
  \end{enumerate}
  \end{tcolorbox}
}

\newenvironment{blocreponses}[1]{%
  \begin{tcolorbox}[enhanced,breakable,colback=bleuclair,colframe=bleufonce,boxrule=0.6pt,arc=2mm,left=2mm,right=2mm,top=2mm,bottom=1.5mm,title={\bfseries #1},coltitle=white,colbacktitle=bleufonce,before skip=3mm,after skip=3mm, fontupper=\large]
  \large
  \begin{enumerate}[leftmargin=*,label=\textbf{\arabic*.},itemsep=0.42em,topsep=0.5em]
  \setcounter{enumi}{\value{reponse}}
}{%
  \setcounter{reponse}{\value{enumi}}
  \end{enumerate}
  \end{tcolorbox}
}

\newtcolorbox{correctionbox}[1]{enhanced,breakable,colback=grisclair,colframe=black!55,boxrule=0.55pt,arc=2mm,left=2mm,right=2mm,top=2mm,bottom=1.5mm,title=\textbf{#1},coltitle=black,colbacktitle=black!12,before skip=2mm,after skip=2mm}
\newtcolorbox{methodebox}[1]{enhanced,breakable,colback=white,colframe=bleu!70!black,boxrule=0.55pt,arc=2mm,left=2mm,right=2mm,top=2mm,bottom=1.5mm,title=\textbf{#1},coltitle=white,colbacktitle=bleu!70!black,before skip=2mm,after skip=2mm}

\newcommand{\titreSujet}[2]{%
  \subsection*{#1}
  \addcontentsline{toc}{subsection}{#1}
  \textit{Référence / inspiration : #2}\par\medskip\hrule\medskip
}
\newcommand{\qcm}[4]{%
\begin{center}\begin{tabularx}{0.96\linewidth}{@{}XX@{}}
\textbf{a.} #1 & \textbf{b.} #2\\[1mm]
\textbf{c.} #3 & \textbf{d.} #4
\end{tabularx}\end{center}}

\tcbset{enhanced,arc=2mm,outer arc=2mm,boxrule=0.42pt,left=1.15mm,right=1.15mm,top=0.75mm,bottom=0.75mm,boxsep=0.35mm,before skip=0.8mm,after skip=0.8mm,fonttitle=\bfseries\footnotesize,coltitle=encre,before upper=\raggedright\fontsize{7.15}{7.65}\selectfont}
\newtcolorbox{fichebox}[3][]{colback=white,colframe=#2!72!black,colbacktitle=#2!18,coltitle=#2!50!black,borderline west={0.9mm}{0pt}{#2!78!black},title={#3},drop fuzzy shadow=black!5,#1}
\newcommand{\tagcol}[2]{\begin{tcolorbox}[enhanced,colback=#1!11,colframe=#1!45!black,boxrule=0.42pt,arc=2mm,left=1mm,right=1mm,top=0.45mm,bottom=0.45mm,before skip=0mm,after skip=0.85mm,fontupper=\bfseries\scriptsize,colupper=#1!55!black]\centering #2\end{tcolorbox}}
\newtcolorbox{panel}[1]{colback=fondpage,colframe=fondpage,boxrule=0pt,arc=0mm,left=0mm,right=0mm,top=0mm,bottom=0mm,height=168mm,valign=top,before skip=0mm,after skip=0mm}

\newcommand{\loitableau}{%
\begin{center}
\begin{tabular}{c|cccc}
$x_i$ & $x_1$ & $x_2$ & $\cdots$ & $x_k$ \\ \hline
$P(X=x_i)$ & $p_1$ & $p_2$ & $\cdots$ & $p_k$
\end{tabular}
\end{center}}

\newcommand{\schemaMoyenne}{%
\begin{center}
\begin{tikzpicture}[scale=0.62,>=Latex]
\draw[->,thick,gris] (-2.4,0)--(2.6,0);
\foreach \x/\lab in {-1.7/$x_1$, -.4/$x_2$, 1.3/$x_3$}{\draw[thick] (\x,.08)--(\x,-.08) node[below] {\scriptsize \lab};}
\draw[very thick,rouge!80!black] (.25,.28)--(.25,-.28) node[above=3pt] {\scriptsize $E(X)$};
\end{tikzpicture}
\end{center}}

\newcommand{\titrepage}{\begin{tcolorbox}[colback=white,colframe=bleu!70!black,boxrule=0.65pt,arc=3mm,left=2.5mm,right=2.5mm,top=1.15mm,bottom=1.15mm,before skip=0mm,after skip=1.4mm,drop fuzzy shadow=black!10]
{\Large\bfseries Parcours de révision -- Variables aléatoires}\hfill {\normalsize Première spécialité -- sans calculatrice}
\end{tcolorbox}}

\newcommand{\ficheRecap}{%
\begin{landscape}
\pagecolor{fondpage}
\thispagestyle{empty}
\titrepage
\begin{tabularx}{\linewidth}{@{}X@{\hspace{1.8mm}}X@{\hspace{1.8mm}}X@{}}
\begin{panel}{}
\tagcol{bleu}{Modéliser}
\begin{fichebox}{bleu}{Variable aléatoire}
Une variable aléatoire $X$ associe un nombre réel à chaque issue d'une expérience aléatoire.\\
Exemple : $X$ peut représenter un gain, un prix payé, un nombre de points, etc.
\end{fichebox}
\begin{fichebox}{bleu}{Valeurs prises}
On commence par lister toutes les valeurs possibles de $X$ :
\[
X(\Omega)=\{x_1;x_2;\ldots;x_k\}.
\]
Il ne faut pas oublier les valeurs négatives quand $X$ est un gain algébrique.
\end{fichebox}
\begin{fichebox}{bleu}{Loi de probabilité}
La loi de $X$ donne toutes les probabilités $P(X=x_i)$.
\loitableau
On doit toujours vérifier :
\[
p_1+p_2+\cdots+p_k=1.
\]
\end{fichebox}
\begin{fichebox}{bleu}{Gain algébrique}
Dans un jeu :
\[
\text{gain algébrique}=\text{somme reçue}-\text{prix payé}.
\]
C'est ce gain que l'on utilise pour décider si le jeu est favorable.
\end{fichebox}
\end{panel}
&
\begin{panel}{}
\tagcol{vert}{Espérance}
\begin{fichebox}{vert}{Formule}
Si $X$ prend les valeurs $x_1,\ldots,x_k$ avec les probabilités $p_1,\ldots,p_k$, alors :
\[
E(X)=p_1x_1+p_2x_2+\cdots+p_kx_k.
\]
\end{fichebox}
\begin{fichebox}{vert}{Interprétation}
$E(X)$ est la valeur moyenne attendue si l'on répète l'expérience un grand nombre de fois.
\end{fichebox}
\begin{fichebox}{vert}{Jeu équitable}
On note $X$ le \textbf{gain algébrique} du joueur.\\[0.3em]
\begin{center}
\renewcommand{\arraystretch}{1.15}
\begin{tabular}{c|c}
$E(X)>0$ & favorable au joueur\\
$E(X)=0$ & équitable\\
$E(X)<0$ & défavorable au joueur
\end{tabular}
\end{center}
\end{fichebox}
\begin{fichebox}{vert}{Phrase type}
\textit{Sur un grand nombre de parties, le gain moyen est d'environ $E(X)$ euro par partie.}
\end{fichebox}
\end{panel}
&
\begin{panel}{}
\tagcol{orange}{Variance et écart type}
\begin{fichebox}{orange}{Variance}
La variance mesure la dispersion autour de l'espérance.\\
Si $m=E(X)$, alors :
\[
V(X)=p_1(x_1-m)^2+p_2(x_2-m)^2+\cdots+p_k(x_k-m)^2.
\]
\end{fichebox}
\begin{fichebox}{orange}{Écart type}
\[
\sigma(X)=\sqrt{V(X)}.
\]
Plus l'écart type est grand, plus les valeurs de $X$ sont dispersées.
\end{fichebox}
\begin{fichebox}{orange}{Comparer deux jeux}
\begin{itemize}[leftmargin=*,itemsep=0.15em]
\item L'espérance compare le gain moyen.
\item La variance ou l'écart type compare le risque.
\item À espérance égale, le jeu de plus grand écart type est le plus risqué.
\end{itemize}
\end{fichebox}
\end{panel}
\end{tabularx}
\clearpage
\nopagecolor
\end{landscape}}

\begin{document}
\begin{center}
{\LARGE\bfseries Corrigé -- Parcours de révision -- Variables aléatoires}\\[1mm]
{\large Première générale -- spécialité mathématiques}\\[1mm]
{\normalsize Sans calculatrice}
\end{center}
\tableofcontents\newpage
\phantomsection\addcontentsline{toc}{section}{Fiche récapitulative}
\ficheRecap

\section{Correction des questions flash}
\setcounter{reponse}{0}
\begin{blocreponses}{Réponses 1 à 12}
\item Oui, car $\frac12+\frac14+\frac14=1$.
\item $P(X=4)=\frac14$.
\item $P(X\geq 1)=\frac14+\frac14=\frac12$.
\item $E(X)=0\times\frac12+1\times\frac14+4\times\frac14=\frac54$.
\item $E(Y)=-2\times\frac35+3\times\frac25=0$.
\item $E(X)=2\times\frac14+8\times\frac34=\frac12+6=\frac{13}{2}$.
\item $V(X)=\frac12(0-2)^2+\frac12(4-2)^2$.
\item $\sigma(X)=\sqrt9=3$.
\item Le jeu est équitable.
\item Non, il est défavorable au joueur.
\item $E(X)=\sum p_ix_i$.
\item $V(X)=p_1(x_1-m)^2+p_2(x_2-m)^2+\cdots+p_k(x_k-m)^2$, avec $m=E(X)$.
\end{blocreponses}
\begin{blocreponses}{Réponses 13 à 24}
\item $P(X=1)=\frac12$ et $P(X=0)=\frac12$.
\item $E(X)=\frac12$.
\item $P(X=1)=\frac36=\frac12$.
\item $G$ prend les valeurs $4$ et $-1$.
\item $P(G=4)=\frac25$ et $P(G=-1)=\frac35$.
\item $E(G)=4\times\frac25+(-1)\times\frac35=1$.
\item $X$ prend les valeurs $7$ et $-3$.
\item $E(X)=7\times\frac15+(-3)\times\frac45=-1$.
\item $E(X)=0\times\frac14+2\times\frac14+5\times\frac12=3$.
\item $V(X)=\frac14(0-3)^2+\frac14(2-3)^2+\frac12(5-3)^2$.
\item $V(X)=\frac94+\frac14+2=\frac92$.
\item Celui qui a le plus grand écart type est le plus risqué.
\end{blocreponses}
\begin{blocreponses}{Réponses 25 à 36}
\item $E(X)=3$.
\item $V(X)=\frac12(1-3)^2+\frac12(5-3)^2=4$.
\item La remise moyenne vaut $5$ euros.
\item Le prix moyen payé vaut $40-6=34$ euros.
\item Il représente souvent le prix payé pour participer.
\item $E(X)=2\times\frac12+6\times\frac12=4$.
\item $V(X)=\sigma(X)^2=4$.
\item $X$ est constante.
\item $E(X)=-1\times\frac14+0\times\frac12+2\times\frac14=\frac14$.
\item $V(X)=\frac14\left(-1-\frac14\right)^2+\frac12\left(0-\frac14\right)^2+\frac14\left(2-\frac14\right)^2$.
\item $V(X)=\frac{25}{64}+\frac{1}{32}+\frac{49}{64}=\frac{19}{16}$.
\item L'écart type mesure la dispersion des valeurs autour de l'espérance.
\end{blocreponses}

\section{Correction du QCM}
\begin{correctionbox}{Réponses du QCM}
\begin{center}
\begin{tabular}{c|cccccccc}
Question & 1&2&3&4&5&6&7&8\\\hline
Réponse & b&b&b&b&b&a&b&c
\end{tabular}
\end{center}
\end{correctionbox}

\section{Correction des sujets type bac / E3C}

\titreSujet{Sujet 1 -- Jeu de jetons et espérance}{correction}
\begin{correctionbox}{1. Valeurs prises par $X$}
Le jeu coûte 2 euros.
\begin{itemize}
\item Si le jeton est rouge, le joueur reçoit 8 euros, donc $X=8-2=6$.
\item Si le jeton est bleu, le joueur reçoit 3 euros, donc $X=3-2=1$.
\item Si le jeton est blanc, le joueur ne reçoit rien, donc $X=0-2=-2$.
\end{itemize}
Ainsi, $X$ prend les valeurs $6$, $1$ et $-2$.
\end{correctionbox}
\begin{correctionbox}{2. Loi de probabilité}
Il y a 10 jetons au total. Donc :
\[
P(X=6)=\frac1{10},\qquad P(X=1)=\frac3{10},\qquad P(X=-2)=\frac6{10}.
\]
\begin{center}
\begin{tabular}{c|ccc}
$x_i$ & $-2$ & $1$ & $6$\\\hline
$P(X=x_i)$ & $\frac6{10}$ & $\frac3{10}$ & $\frac1{10}$
\end{tabular}
\end{center}
\end{correctionbox}
\begin{correctionbox}{3. Espérance}
\[
E(X)=(-2)\times\frac6{10}+1\times\frac3{10}+6\times\frac1{10}.
\]
Donc :
\[
E(X)=\frac{-12+3+6}{10}=-\frac3{10}.
\]
\end{correctionbox}
\begin{correctionbox}{4 à 6. Interprétation et jeu équitable}
Sur un grand nombre de parties, le joueur perd en moyenne $\frac3{10}$ euro par partie. Le jeu est donc défavorable au joueur.

Pour que le jeu soit équitable, le prix d'entrée doit être égal à la somme moyenne reçue. Cette somme moyenne vaut :
\[
8\times\frac1{10}+3\times\frac3{10}+0\times\frac6{10}=\frac{17}{10}.
\]
Il faudrait donc fixer un prix d'entrée de $\frac{17}{10}$ euro, c'est-à-dire 1,70 euro.
\end{correctionbox}

\titreSujet{Sujet 2 -- Deux jeux à comparer}{correction}
\begin{correctionbox}{1 et 2. Lois de probabilité}
Pour le jeu A :
\[
P(A=9)=\frac16,\qquad P(A=3)=\frac26=\frac13,
\qquad P(A=-3)=\frac36=\frac12.
\]
\begin{center}
\begin{tabular}{c|ccc}
$a_i$ & $-3$ & $3$ & $9$\\\hline
$P(A=a_i)$ & $\frac12$ & $\frac13$ & $\frac16$
\end{tabular}
\end{center}
Pour le jeu B :
\[
P(B=3)=\frac12,
\qquad P(B=-3)=\frac12.
\]
\begin{center}
\begin{tabular}{c|cc}
$b_i$ & $-3$ & $3$\\\hline
$P(B=b_i)$ & $\frac12$ & $\frac12$
\end{tabular}
\end{center}
\end{correctionbox}
\begin{correctionbox}{3. Espérances}
\[
E(A)=(-3)\times\frac12+3\times\frac13+9\times\frac16.
\]
Donc :
\[
E(A)=-\frac32+1+\frac32=1.
\]
De plus :
\[
E(B)=(-3)\times\frac12+3\times\frac12=0.
\]
\end{correctionbox}
\begin{correctionbox}{4 et 5. Variances et écarts types}
Pour le jeu A, comme $E(A)=1$, on calcule directement :
\[
V(A)=\frac12(-3-1)^2+\frac13(3-1)^2+\frac16(9-1)^2.
\]
Donc :
\[
V(A)=8+\frac43+\frac{64}{6}=20.
\]
Ainsi :
\[
\sigma(A)=\sqrt{20}=2\sqrt5.
\]
Pour le jeu B, comme $E(B)=0$, on calcule directement :
\[
V(B)=\frac12(-3-0)^2+\frac12(3-0)^2=9.
\]
Donc :
\[
\sigma(B)=3.
\]
\end{correctionbox}
\begin{correctionbox}{6 et 7. Comparaison}
Les deux jeux n'ont pas le même gain moyen :
\[
E(A)=1\qquad \text{et}\qquad E(B)=0.
\]
Le jeu A est plus favorable en moyenne. En revanche :
\[
V(A)=20>9=V(B).
\]
Le jeu A est donc aussi plus risqué, car ses gains sont plus dispersés.
\end{correctionbox}

\titreSujet{Sujet 3 -- Remise commerciale et prix payé}{correction}
\begin{correctionbox}{1 et 2. Valeurs et loi de $P$}
L'article coûte 50 euros.
\begin{itemize}
\item Avec une remise de 20 euros, le client paie $30$ euros.
\item Avec une remise de 10 euros, le client paie $40$ euros.
\item Sans remise, le client paie $50$ euros.
\end{itemize}
Ainsi $P$ prend les valeurs $30$, $40$ et $50$.

Il y a 20 cartes au total :
\[
P(P=30)=\frac2{20}=\frac1{10},\quad
P(P=40)=\frac5{20}=\frac14,
\quad P(P=50)=\frac{13}{20}.
\]
\begin{center}
\begin{tabular}{c|ccc}
$p_i$ & $30$ & $40$ & $50$\\\hline
$P(P=p_i)$ & $\frac1{10}$ & $\frac14$ & $\frac{13}{20}$
\end{tabular}
\end{center}
\end{correctionbox}
\begin{correctionbox}{3 et 4. Espérance}
\[
E(P)=30\times\frac1{10}+40\times\frac14+50\times\frac{13}{20}.
\]
Donc :
\[
E(P)=3+10+\frac{130}{4}=\frac{91}{2}.
\]
Le prix moyen payé est donc $\frac{91}{2}$ euros, c'est-à-dire 45,50 euros.
\end{correctionbox}
\begin{correctionbox}{5 et 6. Variance et écart type}
Comme $E(P)=\frac{91}{2}$, on calcule directement :
\[
V(P)=\frac1{10}\left(30-\frac{91}{2}\right)^2+\frac14\left(40-\frac{91}{2}\right)^2+\frac{13}{20}\left(50-\frac{91}{2}\right)^2.
\]
Donc :
\[
V(P)=\frac{961}{40}+\frac{121}{16}+\frac{1053}{80}=\frac{179}{4}.
\]
Ainsi :
\[
\sigma(P)=\sqrt{\frac{179}{4}}=\frac{\sqrt{179}}{2}.
\]
\end{correctionbox}
\begin{correctionbox}{7. Affirmation du magasin}
La remise moyenne vaut :
\[
50-E(P)=50-\frac{91}{2}=\frac9{2}.
\]
La remise moyenne est donc de 4,50 euros, et non de 5 euros.

L'affirmation du magasin est donc fausse.
\end{correctionbox}

\end{document}
